\honest{Think Like Reality}{Eliezer Yudkowsky}{2007}

When I hear someone call quantum physics ``weird'', bewailing the effects of observation, nonlocal correlations or the impossibility of knowing position and momentum together, I think: ``This person will never understand physics no matter how many books they read.'' \nb{Feynman called nature ``absurd from the point of view of common sense'' in \textsc{QED} (1985). He also gave this post's advice: if you do not like how nature is, ``that's going to get in the way of your understanding it.''}

Reality was here first. ``The universe was propagating complex amplitudes through configuration space for ten billion years before life ever emerged on Earth.'' Quantum physics is not weird. You are weird, for expecting little billiard balls, when ``in fact reality is a perfectly normal cloud of complex amplitude in configuration space.'' \nb{That is one interpretation of quantum mechanics, which takes the wave function to be physically real. Physicists and philosophers have not agreed on it; on epistemic readings the quantum state is not an element of reality. The book argues for its reading of quantum mechanics only in a later sequence.}

Evolution, ``a hack'' that built the retina backward, gave you intuitions for solid objects bouncing around in three dimensions, which is what it took to chase tigers. But tigers are leaky surface generalizations: they arose gradually and are not all alike. At the fundamental level, where laws are stable, global and without exception, there are no tigers, and no such objects either. ``Deal with it.'' \nb{In pilot-wave theory, particles do have positions in ordinary space.}

Surprise is the measure of a poor model. A model that keeps hitting events it called very unlikely should be discarded. A good model makes reality look normal. Intuition is a model too. Poor intuitions are shocked by reality, and good ones make it feel natural. So reshape your intuitions until the universe looks normal: ``think like reality''. You cannot force this, and pretending to feel at home is ``denying your confusion''. But thinking ``How bizarre!'' wastes time, keeps you in the old frame and ``feeds your sense of righteous indignation''. \nb{No evidence is given for these claims about the mind.}

It is the same outside physics. If you ``just don't understand'' how a PhD physicist can believe in astrology, your model of people is at fault: ``People compartmentalize, enough said.'' I now avoid that idiom myself.

``Surprise exists in the map, not in the territory.'' When a fact seems bizarre, first check that it is a fact. If it is, ``the problem isn't the fact, it's you.''
